\section[Cocientes conjugados y memoria mínima]{Cocientes conjugados y
memoria mínima del atlas finito}
\label{sec:cocientes-conjugados-rev11}
\index{cociente!conjugado}
\index{memoria!mínima}
\index{selector interno}
\index{palabra dodecafásica}

La doble proyección exige distinguir el estado transportado de cada una de
sus imágenes. En el nivel finito del atlas, esa distinción admite una
formulación exacta: una misma palabra dodecafásica posee varios cocientes,
cada uno de los cuales conserva una familia determinada de invariantes y
contrae el resto. La configuración denominada históricamente
\emph{diamante} se formula aquí como un sistema de cocientes conjugados sobre
un soporte común. Su contenido matemático consiste en identificar las fibras
de esos cocientes y en determinar la memoria mínima necesaria para separar
los estados que una proyección ha identificado.

\subsection{Dominio celular, rutas y estado interno enriquecido}

Sea
\[
 \mathcal C_9=\{1,\ldots,9\}^2
\]
la carta finita de APP. Para una celda \(x=(r,c)\), el atlas fija el tipo
celular interno, el nivel
\[
 G(r,c)=
 \begin{cases}
 0,&(r,c)=(5,5),\\
 \max\{1,|r-5|,|c-5|\},&\text{en otro caso},
 \end{cases}
\]
y, en el sector correspondiente, la clase de color \((r-c)\bmod 3\). La
evaluación multiplicativa publicada es
\[
 a_{\times}(r,c)=\dr(rc).
\]
Estos lectores se calculan desde la celda orientada; una denominación física
posterior no interviene en su definición.

El catálogo de rutas asocia a \(x\) la coordenada
\[
 \mathcal R(x)=
 \bigl(n_A,s_C,k_{\Delta_4},\nu_{120},\nu_{270},\tau_{12}\bigr)
 \in\mathscr R_{56},
\]
donde \(\tau_{12}\in\{-1,0,1\}^{12}\) es la palabra dodecafásica. La
comparación exhaustiva del atlas de \(81\) celdas con el catálogo CT--108
ampliado establece un único mapa sobreyectivo
\[
 \mathcal R:\mathcal C_9\longrightarrow\mathscr R_{56},
 \qquad |\operatorname{im}\mathcal R|=56,
\]
cuyas fibras tienen la distribución
\begin{equation}
 41\cdot1+10\cdot2+2\cdot3+1\cdot4+2\cdot5=81.
 \label{eq:fibras-r56-rev11}
\end{equation}

Para \(\tau=(\tau_0,\ldots,\tau_{11})\) y
\(d\in\{3,4,6\}\), sean \(\mathcal O_d\) las órbitas de
\(j\mapsto j+d\) en \(\mathbb Z/12\mathbb Z\). Definimos
\[
 \nu_d(\tau)=
 \sum_{O\in\mathcal O_d}
 \operatorname{sgn}\!\left(\sum_{j\in O}\tau_j\right),
 \qquad
 (\nu_{90},\nu_{120},\nu_{180})=(\nu_3,\nu_4,\nu_6),
\]
y las correlaciones cíclicas
\[
 Q_k(\tau)=\sum_{j=0}^{11}\tau_j\tau_{j+k},
 \qquad k=1,\ldots,6,
\]
con índices tomados módulo doce. Completan el lector
\[
 z_0(\tau)=\#\{j:\tau_j=0\},\qquad
 w(\tau)=\sum_j|\tau_j|,\qquad
 L(\tau)=\sum_j\tau_j.
\]
La notación \(z_0\) separa el número de entradas nulas del invariante distinto
\(R_3=Q_3/2\).

\HMTClaim{HMT-R-ATLAS-MINIMAL-INTERNAL-SELECTOR}
\begin{theorem}[Selector interno mínimo]
\label{thm:selector-interno-minimo-rev11}
Fijado el origen de la carta y su orden lexicográfico, sea
\[
 \kappa(x)=
 \#\{y\in\mathcal R^{-1}(\mathcal R(x)):y<_{\rm lex}x\}
 \in M_5:=\{0,1,2,3,4\}.
\]
La aplicación
\begin{align}
 \mathcal S_{\rm int}:\mathcal C_9&\longrightarrow
 \mathscr E_{\rm int}:=\mathscr R_{56}\times M_5\times\mathbb Z^9,
 \notag\\
 x&\longmapsto
 \bigl(\mathcal R(x),\kappa(x),\nu_{90}(\tau_x),
 \nu_{180}(\tau_x),Q_1(\tau_x),\ldots,Q_6(\tau_x),z_0(\tau_x)\bigr)
 \label{eq:selector-interno-rev11}
\end{align}
es inyectiva. Entre los levantamientos inyectivos de la forma
\((\mathcal R,m)\), su alfabeto de memoria tiene cardinal mínimo.
\end{theorem}

\begin{proof}
La distribución \eqref{eq:fibras-r56-rev11} contiene fibras de cardinal
cinco. Por el principio del palomar, todo mapa \(m:\mathcal C_9\to M\) para
el cual \((\mathcal R,m)\) sea inyectivo satisface \(|M|\geq5\). Dentro de
cada fibra, \(\kappa\) enumera exactamente sus elementos de acuerdo con el
orden fijado; por tanto separa dos celdas que comparten ruta. Las celdas con
rutas distintas ya se separan por \(\mathcal R\). En consecuencia,
\((\mathcal R,\kappa)\), y con mayor razón
\(\mathcal S_{\rm int}\), es inyectiva y alcanza la cota inferior.
\end{proof}

La proyección sobre la primera coordenada produce el diagrama conmutativo
\[
\begin{CD}
 \mathcal C_9 @>{\mathcal S_{\rm int}}>> \mathscr E_{\rm int}\\
 @V{\mathcal R}VV @VV{\operatorname{pr}_{\mathscr R}}V\\
 \mathscr R_{56} @>{\operatorname{id}}>> \mathscr R_{56}.
\end{CD}
\]
Así, \(\kappa\) mide la memoria celular mínima olvidada por
\(\mathcal R\). Su canonicidad es relativa al origen y al orden de la carta:
si éstos se consideran un calibre, el objeto intrínseco es la fibra ordenada
o su órbita. El índice \(\kappa\) y la memoria antihoja
\(A_6=(e_0-e_6,\ldots,e_5-e_{11})\) responden a pérdidas de información de
la misma naturaleza categórica, pero cumplen funciones diferentes:
\(\kappa\) separa celdas de una fibra de \(\mathcal R\), mientras \(A_6\)
participa en la reconstrucción del registro dodecafásico.

\subsection{Tipos gruesos y multisecciones naturales}

El cociente
\[
 q:\mathcal C_9\longrightarrow\mathscr Y_{13},
 \qquad
 q(x)=\bigl(\operatorname{familia}(x),G(x)\bigr),
\]
tiene trece valores. Las cardinalidades de sus fibras son
\[
 1,8,16;\qquad 2,4,6,8;\qquad 2,4,6,8;\qquad 4,12.
\]
La salida determinada por \(y\in\mathscr Y_{13}\) es la multisección finita
\begin{equation}
 \mathfrak S(y)=
 \{\mathcal S_{\rm int}(x):q(x)=y\}
 \in\operatorname{Fin}^{+}(\mathscr E_{\rm int}).
 \label{eq:multiseccion-natural-rev11}
\end{equation}

Sea \(\rho(r,c)=(10-r,10-c)\). Las trece fibras son estables bajo
\(\rho\), pero sólo \(q^{-1}(q(5,5))\) contiene un punto fijo. Por ello, en
las doce fibras restantes, una sección puntual \(\rho\)-equivariante no
existe: la multisección \eqref{eq:multiseccion-natural-rev11} es la salida
natural del tipo grueso. La elección de un representante requiere un dato
adicional de orientación o de representación y no se obtiene del nombre
asignado posteriormente al tipo.

\subsection{Conjugación dodecafásica y cocientes de firma}

Sea \(\mathcal T_{12}=\{-1,0,1\}^{12}\) y defínase la involución
\[
 C_M:\mathcal T_{12}\longrightarrow\mathcal T_{12},
 \qquad C_M(\tau)=-\operatorname{rev}(\tau).
\]
Agrupamos los lectores lineales y cuadráticos en
\[
 \mathcal L(\tau)=
 (\nu_{90},\nu_{120},\nu_{180},L)\in\mathbb Z^4,
\]
\[
 \mathcal Q(\tau)=
 (Q_1,\ldots,Q_6,z_0,w)\in\mathbb Z^8.
\]
La conjugación satisface las identidades
\begin{equation}
 \mathcal L(C_M\tau)=-\mathcal L(\tau),
 \qquad
 \mathcal Q(C_M\tau)=\mathcal Q(\tau).
 \label{eq:paridad-conjugacion-rev11}
\end{equation}
Para cualquier lector \(F:\mathcal T_{12}\to I_F\), escribimos
\[
 \tau\sim_F\tau'\quad\Longleftrightarrow\quad F(\tau)=F(\tau'),
 \qquad
 \mathfrak Q_F=\mathcal T_{12}/\!\sim_F.
\]
Cuando \(F\) es equivariante respecto de \(C_M\), la involución desciende a
\(\overline C_{M,F}\) y se obtiene
\[
\begin{CD}
 \mathcal T_{12} @>{C_M}>> \mathcal T_{12}\\
 @V{q_F}VV @VV{q_F}V\\
 \mathfrak Q_F @>{\overline C_{M,F}}>> \mathfrak Q_F.
\end{CD}
\]
En particular, \(\overline C_{M,\mathcal Q}\) es la identidad y
\(\overline C_{M,\mathcal L}\) cambia el signo de la firma publicada. Esta
diferencia de paridad es la estructura conjugada de los dos cocientes; no
son dos evaluaciones de estados iniciales distintos.

El atlas contiene siete formas de palabra dodecafásica con multiplicidades
\begin{equation}
 54,2,2,2,7,7,7.
 \label{eq:multiplicidades-formas-rev11}
\end{equation}
Una forma es fija por \(C_M\); las otras seis constituyen tres pares
conjugados. En cada par, las multiplicidades celulares son \(2\) y \(7\).
En consecuencia, ninguna biyección de \(\mathcal C_9\) puede levantar
\(C_M\) preservando las fibras de forma. La inversión \(\rho\) entrelaza las
palabras en \(49\) celdas y deja \(32\) excepciones. La anti-sección exacta
en este nivel es, por tanto, la correspondencia
\[
 x\longmapsto
 \{y\in\mathcal C_9:\tau_y=C_M(\tau_x)\}.
\]
La igualdad numérica \(54+9+9+9=81\) describe la agrupación de una forma
fija y tres pares de multiplicidad \(2+7\); no constituye una
descomposición en órbitas biyectivas de las celdas.

\subsection{Configuración exacta de cocientes}

Definamos los lectores parciales
\[
 \Sigma_7(\tau)=
 (\nu_{90},\nu_{120},\nu_{180},Q_3,Q_4,Q_6,z_0)
 \in\mathbb Z^7,
\]
\[
 \ell_3(\tau)=(\nu_{90},\nu_{120},\nu_{180})\in\mathbb Z^3,
 \qquad
 p_4(\tau)=(Q_3,Q_4,Q_6,z_0)\in\mathbb Z^4.
\]
Cada lector induce su cociente \(\mathfrak Q_{\Sigma_7}\),
\(\mathfrak Q_{\ell_3}\) o \(\mathfrak Q_{p_4}\). Las tres colisiones
siguientes localizan con precisión qué información retiene cada proyección.

\HMTClaim{HMT-R-ATLAS-CONJUGATE-QUOTIENTS}
\begin{theorem}[Configuración de cocientes conjugados]
\label{thm:configuracion-cocientes-conjugados-rev11}
En el catálogo finito se cumplen simultáneamente las afirmaciones siguientes.
\begin{enumerate}[label=\textup{(\roman*)},leftmargin=2.2em]
 \item Las palabras de borde \(H_\alpha\) y \(H_{\rm APP}\) son distintas
 y pertenecen a órbitas diedrales distintas, pero
 \[
  \Sigma_7(H_\alpha)=\Sigma_7(H_{\rm APP})
  =(-3,-4,-5,3,2,2,6).
 \]
 La coordenada omitida de primer vecino las separa:
 \[
  Q_1(H_\alpha)=4,
  \qquad Q_1(H_{\rm APP})=3.
 \]

 \item El soporte obstructor \(P_\pi\) y la plantilla históricamente
 designada por \(\tau_{\rm lep}\) satisfacen
 \[
  \ell_3(P_\pi)=\ell_3(\tau_{\rm lep})=(-3,-4,-6),
 \]
 mientras
 \[
  p_4(P_\pi)=(7,6,6,3),
  \qquad p_4(\tau_{\rm lep})=(4,3,2,5).
 \]

 \item Las plantillas legadas nominalmente como electrónica y protónica
 publican la misma palabra,
 \[
  \tau_e^{\rm hist}=\tau_p^{\rm hist}=0^{12}.
 \]
 Todo invariante que factorice exclusivamente a través de
 \(\mathcal T_{12}\) toma necesariamente el mismo valor en ambas.
\end{enumerate}
Por tanto, la configuración completa es un abanico de cocientes de un mismo
estado enriquecido,
\begin{equation}
 \mathscr E_{\rm int}
 \longrightarrow
 \left\{
 \begin{array}{c}
  \mathfrak Q_{\mathcal L},\\[1mm]
  \mathfrak Q_{\mathcal Q},\\[1mm]
  \mathfrak Q_{\Sigma_7},\\[1mm]
  \mathfrak Q_{\ell_3},\\[1mm]
  \mathfrak Q_{p_4},\\[1mm]
  \mathcal T_{12},
 \end{array}
 \right.
 \label{eq:abanico-cocientes-rev11}
\end{equation}
y no una correspondencia biyectiva entre constantes matemáticas y clases
físicas.
\end{theorem}

\begin{proof}
Las igualdades de (i)--(iii) resultan de evaluar los lectores definidos
arriba sobre las palabras publicadas. En (i), la desigualdad de \(Q_1\)
demuestra que la coincidencia en \(\mathfrak Q_{\Sigma_7}\) procede de la
coordenada suprimida y no de una identidad de palabras. En (ii), la firma
lineal coincide y el vector par difiere, luego la identificación ocurre sólo
en \(\mathfrak Q_{\ell_3}\). En (iii), la igualdad de palabras implica
\(F(\tau_e^{\rm hist})=F(\tau_p^{\rm hist})\) para toda aplicación
\(F\) cuyo dominio efectivo sea exclusivamente \(\mathcal T_{12}\). Así,
cada igualdad determina una fibra concreta y cada separación identifica la
memoria que debe restaurarse. Esto prueba la interpretación
\eqref{eq:abanico-cocientes-rev11}.
\end{proof}

El teorema ofrece un modelo finito del principio general de doble
proyección: una imagen puede conservar una firma exacta y, a la vez,
identificar estados diferentes. La pareja formada por la firma publicada y
la memoria omitida refina el cociente; la recuperación del estado requiere
reincorporar celda, orientación, sección de banda o la coordenada pertinente
en cada caso.

\subsection{Exactitud de los cocientes conjugados y obstrucciones a la reconstrucción sin memoria}

La construcción parte del atlas finito de \(81\) celdas, sus \(56\)
familias de ruta, las siete formas de sección y la involución \(C_M\), todos
ellos definidos sobre el calendario dodecafásico de \(108\) pasos. Un estado
es la cuádrupla formada por celda, ruta, orientación y sección de banda; esta
cuádrupla precede a cualquier denominación física y fija el dominio de los
cocientes conjugados.

Sobre ese dominio, el índice \(\kappa\) y la minimalidad de cinco estados se
deducen del origen y del orden de la carta. Las igualdades de
\cref{thm:configuracion-cocientes-conjugados-rev11} son igualdades exactas de
firmas: en cada una se especifican el lector conservado y la coordenada de
memoria que el cociente olvida. El recorrido exhaustivo de las \(81\) celdas
reproduce de manera independiente las fibras, los valores de \(\nu_d\), los
cocientes \(Q_k\) y la configuración de diamante; esa enumeración finita
complementa la demostración algebraica y permite reproducirla sin alterar sus
premisas.

La conclusión queda demostrada para estas primitivas explícitas. Dos
prolongaciones requieren mapas adicionales: reconstruir el atlas desde el
núcleo mínimo de APP y construir, para las multisecciones no centrales, una
sección física fina que preserve orientación y representación sin consultar
una masa objetivo. Ninguna de esas prolongaciones interviene en las
igualdades finitas ya establecidas.

La construcción es incompatible con cualquiera de las siguientes
variaciones: un cardinal distinto de \(81\), una imagen distinta de \(56\)
familias, una distribución de fibras distinta de
\eqref{eq:fibras-r56-rev11}, un fallo en la reproducción de
\(\nu_d\) y \(Q_k\), la inyectividad de
\(\mathcal S_{\rm int}\), las tres igualdades del
\cref{thm:configuracion-cocientes-conjugados-rev11}, el recuento
\(49+32\) del entrelazamiento parcial de \(\rho\), o el cierre functorial de
los cobordes del grafo suplementario. Una interpretación física exige,
además, un mapa de representación que conserve esas identidades y un
contraste independiente de su construcción.
